\begin{table}[htbp]
\centering
\footnotesize
\caption{Linearidade cruzada (openai/gpt-4o-mini): fidelidade e acurácia em R$^2$, R$^3$ e R$^4$ nos problemas linear e não lineares; zero-shot do LLM (Passo 2) e métrica em R$^2$/R$^3$/R$^4$ (Passo 4 em A, D e G; Passo 2 no F).}
\label{tab:final_linearidade_cruzada_R2_R3_R4_gpt4mini}
\begin{tabular}{lllrrrrrrrr}
\toprule
 &  &  & \multicolumn{4}{c}{Fidelidade (vs LLM) (\%)} & \multicolumn{4}{c}{Acur.\ vs real (\%)} \\ \cmidrule(lr){4-7} \cmidrule(lr){8-11}
Problema & Atributos & Algoritmo & R$^2$ & R$^3$ & R$^4$ & $\Delta$ R$^4$$-$R$^2$ & R$^2$ & R$^3$ & R$^4$ & $\Delta$ \\
\midrule
A\_linear & x1/x2 & perceptron & 86,7 & 88,7 & 88,9 & +2,2 & 19,3 & 23,6 & 22,9 & +3,6 \\
A\_linear & x1/x2 & nnls & 85,3 & 86,0 & 86,2 & +0,9 & 17,1 & 18,7 & 22,0 & +4,9 \\
\midrule
D\_meialua & x1/x2 & perceptron & — & — & — & — & — & — & — & — \\
D\_meialua & x1/x2 & nnls & — & — & — & — & — & — & — & — \\
\midrule
G1\_homem\_mulher\_classesAB & x1/x2 & perceptron & 73,8 & 74,3 & 75,2 & +1,4 & 18,6 & 19,0 & 18,1 & $-$0,5 \\
G1\_homem\_mulher\_classesAB & x1/x2 & nnls & 74,8 & 73,3 & 82,9 & +8,1 & 20,5 & 21,9 & 27,6 & +7,1 \\
G1\_homem\_mulher\_classesAB & peso/altura & perceptron & 71,4 & 71,4 & 71,4 & +0,0 & 12,4 & 12,4 & 14,3 & +1,9 \\
G1\_homem\_mulher\_classesAB & peso/altura & nnls & 66,2 & 69,5 & 72,9 & +6,7 & 19,5 & 18,1 & 22,4 & +2,9 \\
\midrule
G2\_homem\_mulher & x1/x2 & perceptron & 73,8 & 72,4 & 73,8 & +0,0 & 85,7 & 87,1 & 86,7 & +1,0 \\
G2\_homem\_mulher & x1/x2 & nnls & 78,1 & 73,8 & 74,3 & $-$3,8 & 83,3 & 85,7 & 82,4 & $-$1,0 \\
G2\_homem\_mulher & peso\_kg/altura\_m & perceptron & 95,7 & 96,7 & 95,7 & +0,0 & 91,0 & 91,0 & 91,0 & +0,0 \\
G2\_homem\_mulher & peso\_kg/altura\_m & nnls & 96,7 & 97,1 & 96,7 & +0,0 & 91,0 & 90,5 & 91,0 & +0,0 \\
\bottomrule
\multicolumn{11}{p{0.95\linewidth}}{\footnotesize Síntese cruzada: o ganho de R$^2$ para R$^4$ indica se o LLM adota um critério não linear. A\_linear = Problema A (rótulos zero-shot, Passo 2, métrica do Passo 4); meia-lua = Problema D (150 pontos, Passos 2 e 4); G1, G2 = peso $\times$ altura com split (Passo 2, zero-shot no treino). Fidelidade (vs LLM) e acurácia (vs rótulo real) da métrica, contexto. Uma linha por variante de nomes de atributo (coluna Atributos). Todas as linhas vêm de rótulos zero-shot com classes neutras A/B (A\_linear, meia-lua e as variantes classesAB de G1/G2): a correspondência entre nome e classe é arbitrária e a acurácia (vs rótulo real) está sujeita à permutação em todas elas, não só nas classesAB; apenas em G1/G2 com Mulher/Homem o nome de classe carrega significado. Médias das sementes. ``—'' = sem métrica (o zero-shot do LLM colapsou: classe única ou minoritária com menos de 5 pontos; a linha fica para registrar o colapso). Variantes de G1/G2 cujo zero-shot colapsou em TODAS as sementes não têm linha (o CSV bloco23\_fg\_passos\_2\_e\_4 só grava coletas com métrica; os colapsos estão em bloco3\_g\_passos\_2\_a\_5 e no log).} \\
\end{tabular}
\end{table}
